%!TEX root = ../main.tex
% =============================================================================
%  九、模型检验、灵敏度与稳健性      目标 2.5 页
%  评审在问：你怎么证明是对的？
%  本节含：① 三层验证与 L1--L5 链条（逐项表指向附录 E）
%          ② 非预期性置换检验  ③ 独立审计结论  ④ 基线与消融矩阵（表 27）
%          ⑤ 假设 A9 的量化  ⑥ 灵敏度与结论反转点（表 28）  ⑦ 样本外与伪验证自查
%          ⑧ 口径自洽性检验（三条机器可执行判据）
%  已移入附录 E：L1--L5 五级链条表、A1--A10 逐项审计表、非预期性逐区明细、
%                Sobol 完整指数、容量/功率双向扫描明细、L 形分解对拍、电价过程四项检验
%  红线：求解器返回 optimal 只说明"这个模型被解出来了"，不等于"模型建对了"。
% =============================================================================

\section{模型检验、灵敏度与稳健性}\label{chap:verify}

\subsection{三层验证与审计结论}\label{sec:verify-three}\label{sec:audit}

验证分三层。\textbf{思路层}指假设、信息结构与决策时序，不能靠数值反推，只能由假设与口径声明回答
（\chapref{chap:assump} 的 A1--A11、R1/R2/R3 与 \chapref{chap:q3} 的非预期性约束）。\textbf{模型层}指
变量、约束、目标与单位，“求解器返回最优”\emph{不构成}该层的证明：平衡式漏项或效率方向写反时 LP 同样
返回最优，故本文改用退化恒等式与独立复算替代。\textbf{结果层}指数值、口径与稳健性。三层按 L1--L5
落地：L1 退化与恒等式单测、L2 交叉模型对拍、L3 真实数据独立审计、L4 基线与消融、L5 统计与稳健性；
各级的手段、通过标准与实测值见 \appref{tab:audit-appendix}{E 独立审计报告}（L1--L3）、
表~\ref{tab:ablation}（L4）与表~\ref{tab:sens}（L5）。

\textbf{审计结论。} 审计由独立脚本执行：\textbf{不调用建模的目标函数}，只读原始附件与结果文件重算，
A1--A7 任一失败禁止进入写作、A8--A10 任一失败禁止提交。A1--A10 在内部解（10~min 级）上
\textbf{全部通过}：供给缺口与储电量递推的最大残差分别为 \(3.41\times10^{-13}\) 与
\(1.09\times10^{-11}\)~kWh（阈值 \(10^{-6}\)~kWh），储电量越界数、充放功率越界数、充放互斥违反数与
负购电量数\textbf{均为 0}，费用独立复算与优化目标值的相对偏差 \(\le10^{-6}\)（内部解）、结果文件
回读后 \(4.32\times10^{-7}\)（阈值 \(10^{-4}\)）。交付物层 A3、A4、A6、A7、A9、A10 通过；A1、A5 因
交付物只含逐日 4~h 聚合量而\textbf{不可核验}；A2 的 \(10^{-4}\)~kWh 残差可由“6 个块级充放电量与两个
储电量按 4 位小数取整“的误差上界 \(4\times10^{-4}\)~kWh 完全解释，不是物理越界。覆盖缺口有两处：
内部解审计覆盖问题 1--3 而问题 4 只做交付物层审计；交付物审计按块级口径执行，10~min 级结论只在模型
自身解上成立。

\textbf{最优性证书。} 本文\textbf{不报告“对偶相对间隙”}——多阶段模型的对偶量不构成可行域的最优性
证书。问题 1 改用与求解器无关的证书，即阈值律一致性检验：144 个槽的违反次数为 0、互补松弛最大残差
\(\le10^{-12}\)、\(E_t>0\) 时 \(\alpha_t=p_t\) 的偏差 \(\le10^{-6}\)（适用条件见
命题~\ref{prop:threshold}）。问题 2--4 另用 L 形（Benders）分解这条完全不同的数值路径复核最优值，
结果窗口 334 天上逐日最大相对间隙 \(3.76\times10^{-15}\)、全部收敛（明细见
\appref{sec:lshape}{E L 形分解对拍}）。\evid{C-Q1-02}

\textbf{非预期性置换检验。} 把某一次发布的预报在 12 天之间随机置换后重解，观察四个控制区
（槽 \(0\)--\(36\)、\(36\)--\(72\)、\(72\)--\(108\)、\(108\)--\(144\) 分别由 \(0{:}00\)、\(6{:}00\)、
\(12{:}00\)、\(18{:}00\) 预报控制）的最大变化量。\textbf{任何一次预报被置换，都不改变早于其发布时刻
的决策（变化恰为 0）}，即不存在用未来信息决策的泄漏；而“置换 \(0{:}00\) 预报会改变全天决策”说明检验
本身有效力。\(18{:}00\) 预报的零影响源于其覆盖的傍晚时段光伏实际与预报都接近 0、信息量极低，与
\chapref{chap:q3}“再增加预报时刻最多只能再省 \(58\,062.56\) 元（\(0.42\%\)）”相互印证。
逐区变化量见 \appref{tab:na-app}{E 非预期性置换明细}。\evid{C-Q3-01}

\subsection{基线与消融}\label{sec:ablation}

消融的公平性取三条约束：同一场景集（先固定存盘、所有变体读同一份）、同一指标定义与口径、同一预算。
表~\ref{tab:ablation} 的十组对照除末行外均用同一固定场景集（\(K=60\)、分层 medoid、层样本占比权重）
与同一实际轨迹。

\begin{table}[htbp]
  \centering
  \caption{基线与消融对照矩阵（费用单位：元，电量单位：kWh；全部为 2025-02-01 至 12-31 共 334 天、
           附件 1 固定电价口径；“相对 \(C^{base}\)”以无储能基线 \(16\,407\,319.63\) 元
           （\calref{C-01} 变体 V2）为分母）}
  \label{tab:ablation}
  \scriptsize
  \setlength{\tabcolsep}{3pt}
  \begin{tabularx}{\linewidth}{C{0.72cm} L{3.15cm} C{2.05cm} C{1.55cm} C{1.65cm} Y}
    \toprule
    编号 & 配置（新增或移除的模块） & 总费用 & 相对 \(C^{base}\) & 紧急购电量 & 结论 \\
    \midrule
    B0 & 无储能：净负荷全额购电（闭式 \calref{C-01}）
      & 16\,407\,319.63 & 0 & --- & 基线 \\
    R2 & 典型日预报：计划依据改为附件 1 的日均形状
      & 23\,484\,646.63 & \(-43.14\%\) & 2\,525\,548.76 & 比无储能基线还贵 \(7\,077\,327.00\) 元，不可用 \\
    E0 & 无安全边际：确定性等价（\(K=1\)、零残差情景）
      & 17\,398\,559.68 & \(-6.04\%\) & 1\,156\,618.02 & 移除后反而高于无储能基线；
        \textbf{安全边际的边际价值 \(3\,133\,767.55\) 元} \\
    R1 & 完全信息：\(0{:}00\) 知当日真值
      & 13\,411\,560.90 & 18.26\% & 0.00 & 预测误差成本 \(853\,231.23\) 元（下界） \\
    M & 主模型：\(0{:}00\) 一次预报 \(+\) 两阶段随机规划（\(K=60\)）
      & \textbf{14\,264\,792.13} & 13.06\% & 172\,015.82 & 主方案（基准行） \\
    N2 & 两次预报：\(0{:}00+12{:}00\)
      & 14\,131\,564.16 & 13.87\% & 122\,524.83 & 比主模型省 \(133\,227.98\) 元 \\
    N4 & 四次滚动：\(0/6/12/18\) \(+\) 非预期性
      & 13\,932\,402.95 & 15.08\% & 64\,522.09 & 比主模型省 \(332\,389.18\) 元 \\
    PI & 完全信息上界：四个时刻均知真值
      & 13\,874\,340.39 & 15.44\% & 6\,134.76 & 可获得信息价值的上界；四次预报已实现 \(85.1\%\) \\
    FY & 全年单一 LP：完全信息、单一储电量轨迹、无日决策解耦（192\,384 个变量）
      & 12\,228\,166.56 & 25.47\% & --- & 全局下界；\textbf{日决策解耦代价
        \(1\,183\,394.34\) 元（\(9.68\%\)）} \\
    K120 & 场景数加倍（\(K=120\)，其余同主模型）
      & 14\,235\,163.95 & 13.24\% & 141\,707.94 & 费用仅降 \(29\,628.19\) 元（\(-0.21\%\)），结论稳健 \\
    \bottomrule
  \end{tabularx}
\end{table}

表~\ref{tab:ablation} 中三行需单独说明。① \textbf{安全边际不是装饰}：整行移除后费用升至
\(17\,398\,559.68\) 元、比无储能基线高 \(6.04\%\)，紧急购电量由 \(172\,015.82\) 增至
\(1\,156\,618.02\)~kWh（6.72 倍），故“按情景分位数超购”是问题 2 的核心模块；与实测对照，
逐槽缺货频率 \(8.55\%\)、电量满足率 \(99.04\%\)，落在报童两界 \(1/10\)--\(1/5\) 的\textbf{下侧}
（储能除承担安全边际外还起对冲作用），故报童结论只作方向性近似与上下界 \evid{C-Q2-01}。
② \textbf{场景数不是敏感维度}：\(K\) 由 60 加倍到 120 时费用只降 \(0.21\%\)，但紧急购电量降
\(17.62\%\)、触发天数由 296 降到 277——费用结论稳健，而该分项指标对场景分辨率敏感，引用时
必须带 \(K\)。③ \textbf{日决策解耦的代价}（FY 行）：相差 \(1\,183\,394.34\) 元（\(9.68\%\)），
差额来自“每天固定初态、逐日终端条件”这一工程必需假设，方向保守而非乐观。

\textbf{假设 A9 的量化。} 题面只给光伏预报，未给负载预报，故 A9 把当日负载按已知处理，
此前的声明只是“会低估紧急购电量”。补上量级：只改\textbf{计划阶段}的负载信息，令
\(\widehat L_t=L_t+\varepsilon_t\)，\(\varepsilon\) 为按相对 MAE 标定的零均值误差
（\(\mathrm{MAE}_{\text{rel}}=\mathbb{E}|\varepsilon_t|/\overline L\)，\(\overline L=4626.03\)~kW），
模型、情景机制与结算口径全部不变，重跑 334 天（表~\ref{tab:loaderr}）。

\begin{table}[htbp]
  \centering
  \caption{负载预报误差对问题 2 的影响（只改计划期的负载信息；范围：2025-02-01 至 12-31 共 334 天；
           费用单位元、电量单位 kWh；样本：3 档相对 MAE；\(0\%\) 档与主结果逐位一致）}
  \label{tab:loaderr}
  \small
  \setlength{\tabcolsep}{5pt}
  \begin{tabular}{crrrrc}
    \toprule
    负载预报相对 MAE & 总费用 & 紧急购电量 & 出现天数 & 相对 \(C^{base}\) 节省率 & 相对 \(0\%\) 档费用 \\
    \midrule
    \(0\%\)（主口径） & \(14\,264\,792.13\) & \(172\,015.82\) & \(296\) & \(13.06\%\) & --- \\
    \(5\%\)           & \(15\,611\,126.88\) & \(293\,748.90\) & \(239\) & \(4.85\%\)  & \(+9.44\%\) \\
    \(10\%\)          & \(17\,472\,694.03\) & \(479\,838.20\) & \(202\) & \(-6.49\%\) & \(+22.49\%\) \\
    \bottomrule
  \end{tabular}
\end{table}

两条读数。① \textbf{量级}：相对 MAE 达 \(10\%\) 时问题 2 的总费用上升 \(22.49\%\)，且\textbf{相对无
储能基线的节省率转为负值}（\(-6.49\%\)）——在预报很差的微网里，两阶段计划加紧急购电的组合可能不如
不投运储能。这把 A9 的方向性声明变成可引用的区间：\(13.06\%\) 只适用于负载预报误差可忽略（或与光伏
残差同量级）的场景。② \textbf{一个反直觉但可解释的读数}：误差变大时紧急购电总量上升
（\(172\,015.82\to479\,838.20\)~kWh）而\textbf{触发天数反而下降}（\(296\to202\)）。原因不是计数口径，
而是模型的保守化：零均值误差使情景集变宽，计划层按更宽的分位数多购，“少而大”的缺口取代了“多而小”的
缺口。这也从反面说明主结果里 296 天触发、日均仅 \(581\)~kWh 的形态，是 \(0{:}00\) 光伏预报
\textbf{存在系统性时段偏差}（上午高估、下午低估，见 \ref{sec:sens}）的结果，而非随机噪声的结果。
\evid{C-Q2-01}

\subsection{灵敏度与结论反转点}\label{sec:sens}

灵敏度分两类：\textbf{单参数扫描}回答“参数变多少结论会反转”，\textbf{Sobol 全局灵敏度}回答“哪个参数
决定哪个输出量“。前者见表~\ref{tab:sens}，后者完整指数见 \appref{tab:sobol-app}{E Sobol 灵敏度指数}。

\begin{table}[htbp]
  \centering
  \caption{单参数灵敏度：扰动档位与“结论是否改变”的判断（费用单位 元，电量单位 kWh；
           单参数扫描除终端条件外均在同一组典型日参数下求解；完整档位实测值见
           \appref{tab:sens-app}{E 灵敏度明细}）}
  \label{tab:sens}
  \small
  \setlength{\tabcolsep}{4pt}
  \begin{tabularx}{\linewidth}{L{2.05cm} L{3.15cm} Y}
    \toprule
    参数 & 扰动档位 & 结论是否改变（含反转点） \\
    \midrule
    紧急购电倍数 \(\kappa\) & 2、3、5（基准）、8、10
      & 方向不变：\(\kappa\) 越大安全边际越大、缺口越小。反转点在 \(\kappa\to1\)——
        紧急购电与正常电价同价时缺货成本消失，“需要安全边际”的结论不再成立 \\
    效率口径 \(\eta_c,\eta_d\) & 单边各 0.9（\textbf{主口径}）、往返 0.9 等效单边、
        仅充电侧 0.9、仅放电侧 0.9
      & 区间内结论不变（四种读法下储能均有利可图），但节省额极差 \textbf{13.2\%}
        （式~\eqref{eq:q1-eff}）——同一模型可因效率口径不同报出相差八分之一的数 \\
    储电量运行区间 & 保持 10\%--90\% DOD：容量 7\,200--48\,000~kWh 共 7 档
      & 区间内结论不变：容量摆幅始终是主要瓶颈，仅解除容量上限可再省 22.14\% \\
    充放功率上限 \(\overline P\) & 1\,250、2\,500、5\,000（基准）、7\,500、10\,000~kW
      & 区间内结论不变：功率在 5000~kW 以上边际很小（\(\le0.47\%\)），
        但收紧到 2\,500~kW 时节省率由 26.90\% 降到 24.27\%、1\,250~kW 时降到 16.41\% \\
    终端条件 \calref{C-08} & 窗口：\(S_T\ge S_0\)（主口径）／自由终端；典型日：日周期／自由环状
      & 三组差异均不超过 0.12\%，结论不变；窗口口径下自由终端反而略贵（\(+0.08\%\)），
        说明终端约束并非靠“掏空”储能省钱 \\
    场景数 \(K\) & 60（主口径）、120
      & 费用结论不变（仅降 \(0.21\%\)）；但紧急购电量降 \(17.62\%\)、触发天数由 296 降到 277——
        \textbf{分项指标对场景分辨率敏感，引用时必须带 \(K\)} \\
    \bottomrule
  \end{tabularx}
\end{table}

Sobol 全局灵敏度须\textbf{区分输出量}：对\textbf{典型日全天购电费}起主导作用的是储电量上限
（一阶指数 \(0.2510\)）与下限（\(0.1099\)），其次是往返效率（\(0.0695\)）与充放功率上限
（\(0.0263\)）；对\textbf{节省率}起主导作用的却是电价水平（\(0.8464\)）与负载水平（\(0.4711\)），
储能参数均不超过 \(0.0033\)。原因是问题 1 的确定性模型中，电价、光伏与负载的\emph{水平}对全天购电费
只是整体缩放，不改变调度结构，故一阶指数为 0。结论是“哪一个参数决定哪一个输出量”，而非“所有参数都
重要“。

\subsection{样本外回测与伪验证自查}\label{sec:oos}

场景集与参数在结果窗口的 334 天残差上一次性估计并在同一窗口回放，属\textbf{样本内评估}。故本文
\textbf{不声称样本外有效性}，而以四条代偿性证据约束外推范围：场景保真度（均值相对差 \(0.082\%\)、
关键分位最大相对差 \(6.99\%\)、电价—净需求相关 \(0.2597\) 对 \(0.2620\)）、场景数稳健性
（\(K:60\to120\) 时费用只变 \(-0.21\%\)）、交叉模型对拍（差距方向保守，\(9.68\%\)）与非预期性置换
检验。此外另做两组真正的样本外实验：① 场景集只用 2025-02-01 至 06-30 的残差构造，在测试窗
07-01 至 12-31（184 天）上回放得总费用 \(8\,328\,592.63\) 元，而改用测试窗自身残差构造（样本内上界）
为 \(8\,314\,304.58\) 元，样本外劣化仅 \(14\,288.05\) 元（\(0.172\%\)）；② 四种情景生成器在样本外的
总费用极差仅 \(0.37\%\)（\appref{sec:q4-gen}{F 情景生成对照}），即结论不依赖生成范式，但紧急购电量
的分项差异达 \(48.5\%\)，该分项必须与情景口径同时引用。2025 年 1 月为储电量递推的预热期，不计入结果
（假设 A8）。

\textbf{伪验证自查。} 逐条排除八类伪验证：\textbf{自证循环}（审计脚本不共用建模费用函数）；
\textbf{以求解器状态代替物理校验}（同时报告缺口、递推、越界与互斥残差）；\textbf{比较不同样本}
（四种信息结构、四种预报时刻数、两个场景数与四种生成器均在\emph{同一场景集或同一实际轨迹}上配对
比较）；\textbf{训练集当测试集}（明确声明样本内，另做两组样本外回测）；\textbf{只报点估计}
（关键量均给多档：效率 4 档、容量 7 档、功率 5 档、\(\kappa\) 5 档、\(\lambda\) 11 档、\(\beta\) 6 档、
\(k\) 7 档）；\textbf{事后挑选指标}（全部口径在建模前固定）；\textbf{把“结果好看”当“模型正确”}
（问题 2 的 \(13.06\%\) 小于问题 1 的 \(26.90\%\)，差额可由预测误差成本 \(853\,231.23\) 元解释）；
\textbf{扰动无效}（灵敏度扰动作用于运行区间与功率上限本身，且每档重新求解）。逐条证据见
\appref{sec:fakeverify-app}{J 伪验证自查清单}。

\subsection{口径自洽性检验}\label{sec:caliber-check}

前两节回答“解是否可行、是否最优、是否稳健”，还剩一问：\textbf{同一个数字被报出来时属于哪个口径}。
本题共 11 个口径，而“每天的电价相同”“购电费用”“弃光”在四问之间含义并不相同，口径与数字一旦对不上就
无法复算，故再加三道\textbf{机器可执行}的自洽性判据。

\textbf{检验一（加权均价必须落在电价区间内）。} 对任一含购电费与购电量的结果，令
\(\bar p=\text{报告费用}/\text{报告电量}\)。报告电量是电价向量的逐时段加权平均，故必然满足

\begin{equation}
  \min_t p_t\ \le\ \bar p\ \le\ \max_t p_t ,
  \label{eq:caliber-price-bound}
\end{equation}

否则该费用与该电量不属于同一电价口径。五个主口径的实测值全部通过：问题 1、2、3 的电价向量同为附件 1
（\([0.3713,1.3952]\)~元/kWh），\(\bar p\) 分别为 \(0.5905\)、\(0.6137\)、\(0.5948\)；问题 4-2、4-3
为附件 4（\([0.0076,1.7936]\)），\(\bar p\) 为 \(0.6486\)、\(0.6281\)~元/kWh。\textbf{若某处误用了
全天单一均价}（本题均值 \(0.7662\)），则该口径下所有可行解的 \(\bar p\) 都只能等于 \(0.7662\)，
与上述任何一个值都不符——式~\eqref{eq:caliber-price-bound} 因此能把“报了一个价、算的是另一个价”这类
错配直接暴露。该判据另有一条副产品：问题 2 与问题 3 的\textbf{计划电量完全相同}而计划购电费相差
\(421\,349.41\) 元，恰等于 \calref{C-09} 另一读法下的 \(\sum_tp_tv_t=421\,349.41\) 元
（\S\ref{sec:ablation}），两个口径互相印证。

\textbf{检验二（同一量的三处取值必须一致）。} 每个进入正文的数字都有三个落点：
\texttt{notes/数值台账.md}、\texttt{intermediate/*.json} 与 \texttt{results/result*.xlsx}。做法是从
\textbf{结果文件回读}并重算、再与中间结果对拍：全天购电量的块级最大残差
\(2.96\times10^{-4}\)~kWh（四位小数取整上界 \(7.0\times10^{-4}\)~kWh）、全天购电费最大相对偏差
\(4.32\times10^{-7}\)、日期首末严格为 2025-02-01 与 2025-12-31。题面表 1／表 2 的“指定时段”与槽位的
对应由 \texttt{config.py} 单点定义，正文表格与结果文件因此不可能各写各的。

\textbf{检验三（跨问题比较必须同费用结构）。} 最容易被误用的比较是把“只有计划购电费”与
“计划 \(+\) 紧急 \(+\) 调整相关”两种口径直接相减。本文将信息价值 \calref{C-10} 定义为\textbf{同一
场景集、同一实际轨迹}下仅改变预报时刻数的费用差（\(n=1\) 的 \(14\,264\,792.13\) 元 \(\to\) \(n=4\) 的
\(13\,932\,402.95\) 元），两口径费用项构成相同；问题 4-2 与 4-3 的费用项不同，故只作分项并列、不相减
得节省率。三条判据的合取构成本文的结论防线：\textbf{可行、最优、稳健、口径可分、三处一致、可比者
才比}。\evid{C-Q5-01}
